% 方法節の部品。投稿本文・新規性・結論は未確定。
% P31/P32の凍結仕様に対応。公式クラスでの４ページ組版は未確認。
\section{引き継ぎ状態の比較方法}
\subsection{対象と介入}
モデルの会話履歴を$H$，Pythonの作業状態を$R$，
操作先の模擬環境を$E$，公開済みAPIスキーマを$A$とする．
会話を短縮しても，取得した値が$R$に残る場合がある．
一方，$E$に適用済みの変更と，その変更を行った変数は別の状態である．
本比較は，固定した途中地点で$H$と$R$を操作し，継続への影響を調べる．
担当数や作業の所有権の変更は対象に含めない．

既知の開発用４課題について，元の実行の第４応答後を共通の分岐点とする．
履歴は全ターン保持と直近１完全ターンの保持，
Pythonは作業状態保持と一回初期化の各２条件を交差する．
履歴の短縮は引き継ぎ時に一度だけ行い，以後のターンは両条件で追加する．
初期化条件でも，引き継ぎ後に生成された変数は後続呼び出しへ保持する．
課題要求，公開スキーマ，保存済み観測へのアクセス手段を共通にし，
変数を保持または初期化した事実をモデルに正確に通知する．
したがって，通知の違いも介入に含まれる．
全履歴とは全ての公開ターンを残す条件であり，内部の推論文は含めない．
長い観測のページ分割と原記録の再読手段は全条件で共通にする．

追加の比較は元の４軌跡の第８応答後で行う．
この地点は第４応答後の保持状態量に基づき，
第４応答後の全分岐結果を得る前に固定した．
不適格な地点を別地点へ置き換えず，二つの地点を独立した課題とは数えない．
モデルの乱数列は各地点内の４条件で共通にするが，地点間では異なる．
このため地点間の差を，分岐時期だけの因果効果とは解釈しない．

\subsection{環境とPython状態の再構成}
各分岐に新しい隔離環境を作り，記録した前半のAPIを再実行する．
公開返却，発見済みスキーマ集合，保存DBの内容が一致した場合のみ継続する．
元の環境へ操作を再送しない．Python状態の保持条件は，
元のプログラムに保存したAPI返却を与えて再生する．
このPython再生自体から実APIは送信しない．
要求の名前・引数・順序，標準出力，例外などを元の記録と照合する．
これは観測可能な前半の対応確認であり，任意のヒープ，乱数状態，
全プロセスや未来の挙動の完全な同一性は保証しない．

\subsection{予算と評価}
引き継ぎ後の上限を16生成，入力・出力合計120,000トークン，480秒とする．
課題のAPI上限は写真整理160件，他の３課題80件であり，
元の前半の送信数$N_{\mathrm{pre}}$も課題の上限へ計上する．
継続中の送信数を$N_{\mathrm{suf}}$とすると，
\begin{equation}
 N_{\mathrm{pre}}+N_{\mathrm{suf}}\leq B_{\mathrm{API}}
\end{equation}
を各条件で共通にする．新しい隔離環境の再構成は実験準備として別途記録し，
運用時に無料で状態を移管できるとは仮定しない．
継続時間と再構成を含む実行時間も別に報告する．
使用量不明の生成は除外せず件数を明記し，既知トークン数を下限として扱う．

最終成果は終了後の独立プロセスで評価し，結果を生成側へ戻さない．
未完了や予算到達を成功としない．各課題$i$の４条件の成否$Y_i$から，
履歴削減の影響が保持状態により変わるかを記述するため，
\begin{align}
 d_i={}&(Y_i^{\mathrm{recent,keep}}-Y_i^{\mathrm{full,keep}})\nonumber\\
      &-(Y_i^{\mathrm{recent,reset}}-Y_i^{\mathrm{full,reset}})
\end{align}
を求める．各条件１回の少数開発比較であるため，
$d_i$を安定した個別効果の推定値とは扱わない．
全４件の成否，生成数，トークン，API，再構成費用を併記し，
成功数の増加を伴わない費用減少だけで品質を維持した効率化とは主張しない．
